% ============================================================================
% APPENDIX XG: LIT-GABAIX - Behavioral Inattention & Sparsity
% ============================================================================
% Category: Literature Integration (LIT-)
% Language: English
% EBF Integration: Bounded rationality, attention, and sparse models
% ============================================================================

\chapter*{Appendix UNMAPPED_XG: Xavier Gabaix Research Integration}
\addcontentsline{toc}{chapter}{Appendix UNMAPPED_XG: LIT-GABAIX - Behavioral Inattention \& Sparsity}
\label{app:gabaix}

% ============================================================================
% HEADER BLOCK
% ============================================================================
\begin{tcolorbox}[colback=blue!5!white, colframe=blue!75!black, title=Appendix Metadata]
\textbf{Code:} XG \\
\textbf{Category:} LIT- (Literature Integration) \\
\textbf{Title:} Behavioral Inattention \& Sparsity \\
\textbf{Version:} 1.0 \\
\textbf{Status:} Complete \\
\textbf{Primary Author Focus:} Xavier Gabaix \\
\textbf{Papers Covered:} 20 \\
\textbf{Dependencies:} CORE-AWARE (AU), CORE-WHEN (V), LIT-LAIBSON (LA), DOMAIN-FINANCE \\
\textbf{EBF Connections:} Attention function A($\cdot$), bounded rationality, sparse utility
\end{tcolorbox}

% ============================================================================
% CROSS-REFERENCE MAP
% ============================================================================
\begin{tcolorbox}[colback=green!5!white, colframe=green!75!black, title=Cross-Reference Map]
\textbf{Outgoing References:}
\begin{itemize}[nosep]
    \item CORE-AWARE (AU): Attention modeling $\rightarrow$ Gabaix's sparsity extends awareness
    \item CORE-WHEN (V): Temporal context $\rightarrow$ myopia and discounting
    \item LIT-LAIBSON (LA): Shrouded attributes collaboration
    \item LIT-THALER (R): Bounded rationality foundations
    \item DOMAIN-FINANCE: Financial market applications
    \item METHOD-STAT: Power law methodology
\end{itemize}

\textbf{Incoming References:}
\begin{itemize}[nosep]
    \item Chapter 11 (Awareness): Gabaix inattention model
    \item Chapter 9 (Context): Sparse context representation
    \item PREDICT-01: Attention-based predictions
\end{itemize}
\end{tcolorbox}

% ============================================================================
% ABSTRACT
% ============================================================================
\section*{Abstract}

Xavier Gabaix has developed one of the most influential frameworks for understanding bounded rationality through the lens of \textbf{behavioral inattention} and \textbf{sparsity}. His work provides a mathematically rigorous foundation for why economic agents systematically ignore or underweight certain information. This appendix integrates Gabaix's 20 key papers into the EBF framework, with particular emphasis on:
\begin{enumerate}
    \item The \textbf{Sparsity Model} of bounded rationality
    \item \textbf{Shrouded Attributes} and consumer myopia (with Laibson)
    \item \textbf{Power Laws} in economics and finance
    \item \textbf{Behavioral Macroeconomics} with bounded rationality
    \item \textbf{Granularity Hypothesis} and aggregate fluctuations
\end{enumerate}

% ============================================================================
% QUICK REFERENCE
% ============================================================================
\begin{tcolorbox}[colback=yellow!10!white, colframe=orange!75!black, title=Quick Reference: Key Gabaix Parameters]
\begin{tabular}{lll}
\textbf{Parameter} & \textbf{Value} & \textbf{Source} \\
\hline
$\theta$ (attention) & 0.3-0.7 & Gabaix (2019) \\
Shrouding effect & 40-60\% underweighting & Gabaix \& Laibson (2006) \\
Power law exponent & $\alpha \approx 1$ (Zipf) & Gabaix (1999) \\
Myopia factor & $m \approx 0.5$ & Gabaix \& Laibson (2017) \\
Granularity share & Top 100 firms = 30\% variance & Gabaix (2011) \\
\end{tabular}
\end{tcolorbox}

% ============================================================================
% FUNDAMENTAL QUESTION
% ============================================================================
% -----------------------------------------------------------------------------
% APPENDIX SCOPE BOX
% -----------------------------------------------------------------------------

\begin{tcolorbox}[
    colback=orange!5!white,
    colframe=orange!75!black,
    title={\textbf{Appendix Scope: Ziel / In-Scope / Out-of-Scope / Constraints}},
    fonttitle=\bfseries
]

\textbf{Ziel (Objective):}
\begin{itemize}[nosep]
    \item How do economic agents allocate limited attention, and what are the equilibrium consequences of systematic inattention for individual behavior and aggregate outcomes?
\end{itemize}

\textbf{In-Scope:}
\begin{itemize}[nosep]
    \item Fundamental Question
    \item Shrouded Attributes and Consumer Myopia
    \item Power Laws in Economics
    \item Behavioral Macroeconomics
    \item Granularity and Aggregate Fluctuations
\end{itemize}

\textbf{Out-of-Scope (delegiert an andere Appendices):}
\begin{itemize}[nosep]
    \item Glossary $\rightarrow$ Appendix G
\end{itemize}

\textbf{Constraints (Anwendungsgrenzen):}
\begin{itemize}[nosep]
    \item [TODO: Define when this appendix does NOT apply]
    \item [TODO: Define required prerequisites]
    \item [TODO: Define known limitations]
\end{itemize}

\textbf{Lieferobjekte (Deliverables):}
\begin{itemize}[nosep]
    \item 1 Table(s)
    \item 11 Equation(s)
    \item Worked Example(s)
\end{itemize}

\end{tcolorbox}


\section{Fundamental Question}

\begin{tcolorbox}[colback=red!5!white, colframe=red!75!black, title=Central Question]
\textbf{How do economic agents allocate limited attention, and what are the equilibrium consequences of systematic inattention for individual behavior and aggregate outcomes?}
\end{tcolorbox}

Gabaix's answer: Agents use \textbf{sparse} mental models that focus attention on the most important variables while setting less important ones to default values. This optimization under cognitive constraints explains numerous economic puzzles.

% ============================================================================
% THEORETICAL FOUNDATIONS
% ============================================================================
\section{Theoretical Foundations: The Sparsity Model}

\subsection{Core Framework}

Gabaix's sparsity model (2014, 2019) provides a unified theory of bounded rationality based on attention allocation:

\begin{definition}[Sparse Max]
An agent solves a \textbf{sparse maximization} problem:
\begin{equation}
\max_{a, \mathbf{m}} u(a, \mathbf{x}^d + \mathbf{m} \odot (\mathbf{x} - \mathbf{x}^d)) - \kappa \sum_i |m_i|
\end{equation}
where:
\begin{itemize}
    \item $\mathbf{x}$ = true state vector
    \item $\mathbf{x}^d$ = default/prior state
    \item $\mathbf{m}$ = attention vector ($m_i \in [0,1]$)
    \item $\kappa$ = cognitive cost per unit of attention
    \item $\odot$ = element-wise multiplication
\end{itemize}
\end{definition}

\subsection{Key Axioms}

\begin{axiom}[UNMAPPED_XG-1: Sparse Attention]
\label{ax:xg1}
Attention is allocated sparsely: agents set $m_i = 0$ for variables whose marginal impact falls below the cognitive cost threshold $\kappa$:
\begin{equation}
m_i = 0 \text{ if } \left| \frac{\partial u}{\partial x_i} \right| < \kappa
\end{equation}
\textbf{EBF Connection:} Extends CORE-AWARE (AU) with explicit attention allocation mechanism.
\end{axiom}

\begin{axiom}[UNMAPPED_XG-2: Default Anchoring]
\label{ax:xg2}
Unattended variables are processed at their \textbf{default values} $x_i^d$, typically the mean or mode of the distribution:
\begin{equation}
\tilde{x}_i = x_i^d + m_i(x_i - x_i^d)
\end{equation}
\textbf{EBF Connection:} Formalizes anchoring bias in terms of attention allocation.
\end{axiom}

\begin{axiom}[UNMAPPED_XG-3: Endogenous Salience]
\label{ax:xg3}
Attention allocation responds to variable importance:
\begin{equation}
m_i^* = \max\left(0, 1 - \frac{\kappa}{\sigma_i \cdot |u_{ai}|}\right)
\end{equation}
where $\sigma_i$ is the variance of $x_i$ and $u_{ai}$ is the cross-partial derivative.
\textbf{EBF Connection:} Links context (V) to attention through salience.
\end{axiom}

% ============================================================================
% SHROUDED ATTRIBUTES
% ============================================================================
\section{Shrouded Attributes and Consumer Myopia}

\subsection{The Shrouding Model}

Gabaix \& Laibson (2006) analyze markets where firms can ``shroud'' product attributes:

\begin{theorem}[Shrouded Equilibrium]
In competitive markets with myopic consumers ($\theta < 1$), shrouding persists in equilibrium:
\begin{enumerate}
    \item Firms shroud high add-on prices (e.g., printer ink, hotel wifi)
    \item Sophisticated consumers exploit shrouding but don't eliminate it
    \item De-shrouding strategies are not profitable when enough consumers remain myopic
\end{enumerate}
\end{theorem}

\subsection{Key Findings}

\begin{tcolorbox}[colback=blue!5!white, colframe=blue!75!black]
\textbf{Shrouded Attributes Predictions (Gabaix \& Laibson 2006):}
\begin{itemize}
    \item Add-on prices can exceed competitive levels by 40-60\%
    \item Equilibrium features cross-subsidization: base products below cost, add-ons above
    \item Consumer education campaigns have limited effectiveness
    \item Mandatory disclosure may backfire by enabling price discrimination
\end{itemize}
\end{tcolorbox}

\subsection{EBF Integration: Ψ-Shrouding}

The shrouding model enriches the EBF context dimension:

\begin{equation}
\Psi_{shroud} = \{visibility_j, complexity_j, timing_j, framing_j\}_{j \in attributes}
\end{equation}

Behavioral response depends on attention-adjusted perception:
\begin{equation}
\tilde{p}_{total} = p_{base} + \theta \cdot p_{addon} \quad \text{where } \theta \in [0.3, 0.7]
\end{equation}

% ============================================================================
% POWER LAWS
% ============================================================================
\section{Power Laws in Economics}

\subsection{Zipf's Law for Cities}

Gabaix (1999) provided the first compelling explanation of Zipf's Law:

\begin{theorem}[Zipf from Gibrat]
If cities grow according to Gibrat's Law (proportional growth with i.i.d. shocks) and there's a reflecting lower bound, the limiting distribution follows Zipf's Law:
\begin{equation}
P(Size > S) \propto S^{-\alpha}, \quad \alpha \approx 1
\end{equation}
\end{theorem}

\subsection{Financial Market Fluctuations}

Gabaix et al. (2003, 2006) extended power law analysis to finance:

\begin{axiom}[UNMAPPED_XG-4: Power Law Returns]
\label{ax:xg4}
Financial returns exhibit power-law tails:
\begin{equation}
P(|r| > x) \propto x^{-\alpha}, \quad \alpha \approx 3
\end{equation}
This emerges from the trading behavior of large institutional investors.
\textbf{EBF Connection:} Explains fat tails in behavioral predictions.
\end{axiom}

% ============================================================================
% BEHAVIORAL MACROECONOMICS
% ============================================================================
\section{Behavioral Macroeconomics}

\subsection{The Behavioral New Keynesian Model}

Gabaix (2020) introduces bounded rationality into standard macro models:

\begin{definition}[Behavioral IS Curve]
\begin{equation}
x_t = M \mathbb{E}_t[x_{t+1}] - \sigma(i_t - \mathbb{E}_t[\pi_{t+1}] - r^n)
\end{equation}
where $M < 1$ is the ``cognitive discount factor'' capturing myopia about future conditions.
\end{definition}

\begin{definition}[Behavioral Phillips Curve]
\begin{equation}
\pi_t = \beta M^f \mathbb{E}_t[\pi_{t+1}] + \kappa x_t
\end{equation}
where $M^f < 1$ captures firms' bounded attention to future inflation.
\end{definition}

\subsection{Key Implications}

\begin{axiom}[UNMAPPED_XG-5: Behavioral Macro Dynamics]
\label{ax:xg5}
With $M < 1$:
\begin{enumerate}
    \item Forward guidance becomes less powerful (``forward guidance puzzle'' solved)
    \item Determinacy holds even at the zero lower bound
    \item Fiscal multipliers are larger
    \item Inflation expectations are anchored more easily
\end{enumerate}
\textbf{EBF Connection:} Provides macro-level implications of CORE-AWARE.
\end{axiom}

% ============================================================================
% GRANULARITY HYPOTHESIS
% ============================================================================
\section{Granularity and Aggregate Fluctuations}

\subsection{The Granular Origins}

Gabaix (2011) shows that firm-level shocks don't wash out in aggregation:

\begin{theorem}[Granular Hypothesis]
In an economy with fat-tailed firm size distribution:
\begin{equation}
\sigma_{GDP} = \frac{\sigma_{firm}}{\sqrt{N}} \cdot h \approx \frac{\sigma_{firm}}{\sqrt{N}} \cdot \frac{1}{\ln N}
\end{equation}
where $h$ is the Herfindahl concentration. With Zipf-distributed firms, this explains 30\% of GDP volatility.
\end{theorem}

\begin{axiom}[UNMAPPED_XG-6: Micro-Macro Bridge]
\label{ax:xg6}
Individual-level behavioral heterogeneity can have aggregate consequences when:
\begin{enumerate}
    \item Size distributions are fat-tailed
    \item Network connections amplify individual shocks
    \item Attention cascades create correlated behavior
\end{enumerate}
\textbf{EBF Connection:} Links individual bounded rationality to system-level outcomes.
\end{axiom}

% ============================================================================
% EBF INTEGRATION
% ============================================================================
\section{EBF Integration}

\subsection{10C Mapping}

\begin{table}[h]
\centering
\begin{tabular}{lll}
\textbf{10C Element} & \textbf{Gabaix Contribution} & \textbf{Parameter} \\
\hline
CORE-WHO (AAA) & Agent heterogeneity in $\theta$ & $\theta \in [0.3, 0.9]$ \\
CORE-WHAT (C) & Sparse utility dimensions & $m_i \in \{0, 1\}$ \\
CORE-HOW (B) & Attention complementarities & $\gamma_{attention}$ \\
CORE-WHEN (V) & Temporal attention decay & $M < 1$ \\
CORE-WHERE (BBB) & Attention parameters & $\kappa, \theta$ \\
CORE-AWARE (AU) & Sparse attention model & $A(\cdot) = m \odot x$ \\
CORE-READY (AV) & Action threshold under sparsity & Lower with $\theta < 1$ \\
CORE-STAGE (AW) & Default-based progression & Anchoring dynamics \\
\end{tabular}
\caption{Gabaix contributions mapped to 10C framework}
\end{table}

\subsection{Key Parameter Values for BBB}

\begin{tcolorbox}[colback=green!5!white, colframe=green!75!black, title=Gabaix Parameters for BBB Repository]
\begin{itemize}
    \item $\theta_{general}$: 0.5 (SD 0.15) - General attention parameter
    \item $\theta_{financial}$: 0.3-0.4 - Financial product attention
    \item $\kappa_{cognitive}$: Context-dependent, calibrated to data
    \item $M_{macro}$: 0.85 - Cognitive discount factor
    \item $\alpha_{power}$: 1.0 (Zipf), 3.0 (returns)
\end{itemize}
\end{tcolorbox}

% ============================================================================
% WORKED EXAMPLE
% ============================================================================
\section{Worked Example: Credit Card Fee Attention}

Consider a consumer choosing a credit card with:
\begin{itemize}
    \item Annual fee: \$95
    \item Late payment fee: \$39
    \item Foreign transaction fee: 3\%
    \item Interest rate: 24.99\% APR
\end{itemize}

\textbf{Sparse Max Analysis:}

With attention parameter $\theta = 0.4$ and defaults at industry averages:
\begin{align}
\tilde{Cost}_{annual} &= 95 \quad \text{(fully attended, salient)} \\
\tilde{Cost}_{late} &= 29 + 0.4(39 - 29) = 33 \quad \text{(partially attended)} \\
\tilde{Cost}_{foreign} &= 0 \quad \text{(ignored, below threshold)} \\
\tilde{Cost}_{interest} &= 18 + 0.3(25 - 18) = 20.1\% \quad \text{(underweighted)}
\end{align}

\textbf{Prediction:} Consumer underestimates total expected cost by 25-40\%, consistent with Gabaix \& Laibson (2006) findings.

% ============================================================================
% SUMMARY
% ============================================================================

% === AUTO-GENERATED ENTRIES (2026-02-22) ===

%%%%%%%%%%%%%%%%%%%%%%%%%%%%%%%%%%%%%%%%%%%%%%%%%%%%%%%%%%%%%%%%%%%%%%%%%%%%%%%
\subsubsection{2002: Institutional Investors and Stock Market Volatility}
%%%%%%%%%%%%%%%%%%%%%%%%%%%%%%%%%%%%%%%%%%%%%%%%%%%%%%%%%%%%%%%%%%%%%%%%%%%%%%%

\paragraph{Citation.}
Gabaix, Xavier and Gopikrishnan, Parameswaran and Plerou, Vasiliki and Stanley, H. Eugene (2002). ``Institutional Investors and Stock Market Volatility.'' \textit{Quarterly Journal of Economics}, 121.

\paragraph{10C Integration.}
\begin{itemize}[nosep]
  \item \textbf{Domain}: [To be specified]
  \item \textbf{use\_for}: LIT-GABAIX
\end{itemize}


%%%%%%%%%%%%%%%%%%%%%%%%%%%%%%%%%%%%%%%%%%%%%%%%%%%%%%%%%%%%%%%%%%%%%%%%%%%%%%%
\subsubsection{2007: Why Has CEO Pay Increased So Much? A Talent-Based Explanatio...}
%%%%%%%%%%%%%%%%%%%%%%%%%%%%%%%%%%%%%%%%%%%%%%%%%%%%%%%%%%%%%%%%%%%%%%%%%%%%%%%

\paragraph{Citation.}
Gabaix, Xavier and Landier, Augustin (2007). ``Why Has CEO Pay Increased So Much? A Talent-Based Explanation.'' \textit{Working Paper}.

\paragraph{10C Integration.}
\begin{itemize}[nosep]
  \item \textbf{Domain}: [To be specified]
  \item \textbf{use\_for}: LIT-GABAIX
\end{itemize}


%%%%%%%%%%%%%%%%%%%%%%%%%%%%%%%%%%%%%%%%%%%%%%%%%%%%%%%%%%%%%%%%%%%%%%%%%%%%%%%
\subsubsection{2014: Housing Booms and City Centers}
%%%%%%%%%%%%%%%%%%%%%%%%%%%%%%%%%%%%%%%%%%%%%%%%%%%%%%%%%%%%%%%%%%%%%%%%%%%%%%%

\paragraph{Citation.}
Gabaix, Xavier and Landier, Augustin and Thesmar, David (2014). ``Housing Booms and City Centers.'' \textit{American Economic Review}, 104.

\paragraph{10C Integration.}
\begin{itemize}[nosep]
  \item \textbf{Domain}: [To be specified]
  \item \textbf{use\_for}: LIT-GABAIX
\end{itemize}


%%%%%%%%%%%%%%%%%%%%%%%%%%%%%%%%%%%%%%%%%%%%%%%%%%%%%%%%%%%%%%%%%%%%%%%%%%%%%%%
\subsubsection{2016: Behavioral Macroeconomics via Sparse Dynamic Programming}
%%%%%%%%%%%%%%%%%%%%%%%%%%%%%%%%%%%%%%%%%%%%%%%%%%%%%%%%%%%%%%%%%%%%%%%%%%%%%%%

\paragraph{Citation.}
Gabaix, Xavier (2016). ``Behavioral Macroeconomics via Sparse Dynamic Programming.'' \textit{NBER Working Paper}.

\paragraph{10C Integration.}
\begin{itemize}[nosep]
  \item \textbf{Domain}: [To be specified]
  \item \textbf{use\_for}: LIT-GABAIX
\end{itemize}


%%%%%%%%%%%%%%%%%%%%%%%%%%%%%%%%%%%%%%%%%%%%%%%%%%%%%%%%%%%%%%%%%%%%%%%%%%%%%%%
\subsubsection{2017: A Behavioral Model of the Asset Price Boom and Bust Cycle}
%%%%%%%%%%%%%%%%%%%%%%%%%%%%%%%%%%%%%%%%%%%%%%%%%%%%%%%%%%%%%%%%%%%%%%%%%%%%%%%

\paragraph{Citation.}
Gabaix, Xavier (2017). ``A Behavioral Model of the Asset Price Boom and Bust Cycle.'' \textit{Working Paper}.

\paragraph{10C Integration.}
\begin{itemize}[nosep]
  \item \textbf{Domain}: [To be specified]
  \item \textbf{use\_for}: LIT-GABAIX
\end{itemize}


%%%%%%%%%%%%%%%%%%%%%%%%%%%%%%%%%%%%%%%%%%%%%%%%%%%%%%%%%%%%%%%%%%%%%%%%%%%%%%%
\subsubsection{2021: In Search of the Origins of Financial Fluctuations: The Inel...}
%%%%%%%%%%%%%%%%%%%%%%%%%%%%%%%%%%%%%%%%%%%%%%%%%%%%%%%%%%%%%%%%%%%%%%%%%%%%%%%

\paragraph{Citation.}
Gabaix, Xavier and Koijen, Ralph (2021). ``In Search of the Origins of Financial Fluctuations: The Inelastic Markets Hypothesis.'' \textit{NBER Working Paper}.

\paragraph{10C Integration.}
\begin{itemize}[nosep]
  \item \textbf{Domain}: [To be specified]
  \item \textbf{use\_for}: LIT-GABAIX
\end{itemize}

% === END AUTO-GENERATED ENTRIES ===
\section{Summary}

\begin{tcolorbox}[colback=gray!10!white, colframe=gray!75!black, title=Key Takeaways]
\begin{enumerate}
    \item \textbf{Sparsity Model:} Bounded rationality emerges from optimal attention allocation under cognitive costs
    \item \textbf{Shrouded Attributes:} Firms exploit consumer inattention; equilibrium shrouding persists
    \item \textbf{Power Laws:} Fat tails in economics arise from growth processes and heterogeneity
    \item \textbf{Behavioral Macro:} Cognitive discount factor $M < 1$ solves forward guidance puzzle
    \item \textbf{Granularity:} Micro heterogeneity has macro consequences with fat-tailed distributions
    \item \textbf{EBF Integration:} Gabaix provides rigorous foundations for CORE-AWARE
\end{enumerate}
\end{tcolorbox}

% ============================================================================
% GLOSSARY
% ============================================================================
\section{Glossary}

See Appendix G (Glossary) for standard EBF terms. Gabaix-specific terms:

\begin{description}
    \item[Behavioral Inattention] Systematic underweighting of information due to cognitive costs
    \item[Cognitive Discount Factor ($M$)] Degree to which agents underweight future variables; $M=1$ is rational, $M<1$ is myopic
    \item[Shrouded Attribute] Product feature that firms obscure and consumers underweight
    \item[Sparse Max] Optimization under cognitive constraints with L1 penalty on attention
    \item[Sparsity] Property of mental models that set unimportant variables to defaults
\end{description}

% ============================================================================
% CRITICAL FOUNDATIONS
% ============================================================================
\section{Critical Foundations}

\begin{enumerate}
    \item Gabaix, X. (2019). Behavioral Inattention. \textit{Handbook of Behavioral Economics}. [Survey]
    \item Gabaix, X. (2014). A Sparsity-Based Model of Bounded Rationality. \textit{QJE}. [Core Theory]
    \item Gabaix, X. \& Laibson, D. (2006). Shrouded Attributes. \textit{QJE}. [Consumer Behavior]
    \item Gabaix, X. (2020). A Behavioral New Keynesian Model. \textit{AER}. [Macro Applications]
    \item Gabaix, X. (2011). The Granular Origins of Aggregate Fluctuations. \textit{Econometrica}. [Aggregation]
\end{enumerate}

% ============================================================================
% OPEN ISSUES
% ============================================================================
\section{Open Issues}

\begin{enumerate}
    \item \textbf{Attention measurement:} How to directly elicit $\theta$ parameters in field settings?
    \item \textbf{Attention dynamics:} How does attention evolve with learning and experience?
    \item \textbf{Collective attention:} How do social networks shape aggregate attention allocation?
    \item \textbf{Optimal disclosure:} What information presentation minimizes inattention distortions?
\end{enumerate}

% ============================================================================
% REFERENCES
% ============================================================================
\section{References}

\nocite{bcm_master}

For complete bibliography, see Master Bibliography (\texttt{bibliography/master\_references.bib}).

Key Gabaix references integrated in this appendix:
\begin{itemize}[nosep]
    \item gabaix2019behavioral, gabaix2014sparsity, gabaix2006shrouded
    \item gabaix2017myopia, gabaix2020behavioral\_nk, gabaix2011granular
    \item gabaix1999zipf, gabaix2003theory, gabaix2008ceo
    \item gabaix2012variable, gabaix2008limits
\end{itemize}

\end{document}
